% Cola de depuracao do sonar modificado (exp6) para consulta rapida na prova.
% Compilar: latexmk -lualatex cola_exp6.tex
\documentclass[9pt,a4paper]{extarticle}

\usepackage{fontspec}
\usepackage[brazilian]{babel}
\usepackage[a4paper,margin=11mm,top=10mm,bottom=12mm]{geometry}
\usepackage{tabularx,booktabs,array}
\usepackage{xcolor}
\usepackage{enumitem}
\usepackage{titlesec}
\usepackage{tikz}
\usetikzlibrary{positioning,calc,arrows.meta,decorations.pathreplacing}

\setmainfont{TeX Gyre Heros}
\setmonofont{DejaVu Sans Mono}[Scale=0.88]

\titleformat{\section}{\normalsize\bfseries\color{blue!55!black}}{}{0pt}{}
\titlespacing*{\section}{0pt}{7pt}{3pt}
\setlength{\parindent}{0pt}
\setlength{\parskip}{3pt}
\setlist{nosep,leftmargin=12pt}
\renewcommand{\arraystretch}{1.12}
\setlength{\tabcolsep}{4pt}
\newcolumntype{Y}{>{\raggedright\arraybackslash}X}
\newcommand{\s}[1]{\texttt{#1}}          % sinal / codigo
\newcommand{\bancada}{\textcolor{red!70!black}{\textbf{(06/10)}}}
\pagestyle{plain}

\begin{document}

{\large\bfseries Cola de depuração --- Sonar exp6 (PCS3645)}\hfill
{\small T1BA5 · top \s{sonar2} · 08/10/2026}

% ==========================================================================
\section{Guia de montagem (DE0-CV)}

Placa desligada até o passo 6. É a montagem da exp5 mais o TXD do CP2102.
Conferir o pino 1 de cada conector pela marcação da placa.

\begin{center}
\begin{tikzpicture}[
    pin/.style={circle, draw=black!35, fill=white, inner sep=0pt, minimum size=2.3mm},
    pwr/.style={pin, draw=red!60!black, fill=red!55},
    gnd/.style={pin, draw=black, fill=black!75},
    sig/.style={pin, draw=blue!60!black, fill=blue!45},
    tag/.style={font=\footnotesize, inner sep=1pt},
    num/.style={font=\tiny, text=black!55},
    ld/.style={draw=black!55}
  ]
  \def\sp{0.36}
  % \conector{nome}{x}{rotulo}: 20 linhas de 2 pinos (impares a esquerda)
  \newcommand{\conector}[3]{
    \draw[rounded corners=1pt, draw=black!35, fill=black!4]
      (#2-0.2, 0.2) rectangle (#2+0.62, -19*\sp-0.2);
    \foreach \l in {0,...,19} {
      \pgfmathtruncatemacro{\pinimpar}{2*\l+1}
      \pgfmathtruncatemacro{\pinpar}{2*\l+2}
      \node[pin] (#1-\pinimpar) at (#2, -\l*\sp) {};
      \node[pin] (#1-\pinpar)   at (#2+0.42, -\l*\sp) {};
    }
    \node[font=\footnotesize\bfseries, anchor=south] at (#2+0.21, 0.25) {#3};
  }
  % rotulos: \esq{pino}{texto} a esquerda (impares), \dir{pino}{texto} a direita (pares)
  \newcommand{\esq}[2]{\draw[ld] (#1.west) -- ++(-0.3,0) node[tag, anchor=east] {#2};}
  \newcommand{\dir}[2]{\draw[ld] (#1.east) -- ++(0.3,0) node[tag, anchor=west] {#2};}

  % JP1 = GPIO_0: serial e servo
  \conector{JP1}{5.6}{JP1 · GPIO\_0}
  \node[gnd] at (JP1-12) {};
  \node[sig] at (JP1-2) {};  \node[sig] at (JP1-4) {};  \node[sig] at (JP1-40) {};
  \node[num, anchor=east] at (JP1-1.west) {1};
  \node[num, anchor=east] at (JP1-39.west) {39};
  \esq{JP1-11}{11 · 5\,V da placa: \textbf{não usar} (servo na fonte externa)}
  \dir{JP1-2}{2 · \s{saida\_serial} $\to$ RXD do CP2102}
  \dir{JP1-4}{4 · \s{entrada\_serial} $\leftarrow$ TXD do CP2102}
  \dir{JP1-12}{12 · GND: $-$ da fonte do servo, servo (marrom), CP2102}
  \dir{JP1-40}{40 · \s{pwm} $\to$ sinal do servo (laranja)}

  % JP2 = GPIO_1: sensor e fim_posicao
  \conector{JP2}{13.4}{JP2 · GPIO\_1}
  \node[pwr] at (JP2-29) {};  \node[gnd] at (JP2-30) {};
  \node[sig] at (JP2-2) {};  \node[sig] at (JP2-4) {};  \node[sig] at (JP2-40) {};
  \node[num, anchor=east] at (JP2-1.west) {1};
  \node[num, anchor=east] at (JP2-39.west) {39};
  \esq{JP2-29}{29 · 3,3\,V $\to$ VCC do HC-SR04 (\textbf{nunca 5\,V})}
  \dir{JP2-2}{2 · \s{trigger} $\to$ TRIG}
  \dir{JP2-4}{4 · \s{echo} $\leftarrow$ ECHO}
  \dir{JP2-30}{30 · GND: HC-SR04, AD}
  \dir{JP2-40}{40 · \s{fim\_posicao} (só DIO1)}
\end{tikzpicture}
\end{center}

\textcolor{red!60!black}{$\bullet$} alimentação ·
$\bullet$ GND (todos ligados entre si) ·
\textcolor{blue!60!black}{$\bullet$} sinal · $\circ$ não ligar.
Ímpares à esquerda, pares à direita; sinais só na coluna par (2, 4, 40). A placa só alimenta o sensor (3,3\,V,
pino 29, vizinho do GND 30: um fio deslocado dá curto). O servo usa a fonte
de 5\,V externa ($+$ no VCC do servo, vermelho; $-$ no GND comum, junto com o
marrom do servo); o 5\,V da placa (pino 11) fica livre.

\textbf{Analog Discovery:} cada fio do AD entra no mesmo ponto do sinal, em
paralelo (coluna da protoboard ou pino livre), sem tirar o fio que já está lá.

\begin{tabularx}{\linewidth}{@{}l l Y l l@{}}
\toprule
Fio do AD & Cor* & Prender em & Sinal (pino FPGA) & Ver com \\
\midrule
1$+$  & laranja        & sinal do servo (JP1 · 40)           & \s{pwm} (T15)            & Scope CH1 \\
2$+$  & azul           & TRIG do HC-SR04 (JP2 · 2)           & \s{trigger} (A12)        & Scope CH2 \\
DIO0  & rosa           & ECHO do HC-SR04 (JP2 · 4)           & \s{echo} (B12)           & Logic \\
DIO1  & rosa           & direto no pino JP2 · 40 (não vai em mais nada) & \s{fim\_posicao} (K16) & Logic \\
DIO6  & rosa           & TXD do CP2102 (JP1 · 4)             & \s{entrada\_serial} (C16) & Logic / Protocol UART \\
DIO7  & rosa           & RXD do CP2102 (JP1 · 2)             & \s{saida\_serial} (B16)  & Logic / Protocol UART \\
1$-$, 2$-$ & laranja/branco, azul/branco & GND comum      & referência dos canais     & --- \\
GND   & preto          & GND comum (JP2 · 30)                & referência do AD          & --- \\
\bottomrule
\end{tabularx}

{\footnotesize * Cores do cabo do Analog Discovery 2; conferir pelo rótulo do
conector do cabo. Configurações de cada instrumento na seção Analog Discovery.}

\begin{minipage}[t]{0.49\linewidth}
\vspace{0pt}
\textbf{Montagem e energização}
\begin{enumerate}
  \item \textbf{GND primeiro:} GND da placa (pinos 12 e 30) no GND do sensor,
        do CP2102, do AD e no $-$ da fonte externa do servo.
  \item \textbf{Alimentação:} HC-SR04 em 3,3\,V (JP2 · 29); em 5\,V o echo
        chega a 5\,V na FPGA. Servo na fonte externa de 5\,V ($+$ no
        vermelho, $-$ no marrom); nunca pelo 5\,V da placa nem por pino de E/S.
  \item \textbf{Sinais:} conforme o diagrama. TXD e RXD cruzados; fios do AD (tabela)
        em paralelo, sem interromper o fio.
  \item \textbf{Mecânica:} sensor sobre o eixo do servo, girando livre de
        20° a 160° sem puxar cabos; 90° apontando para a frente.
  \item \textbf{Multímetro, placa desligada:} sem curto entre 3,3\,V e GND, nem entre $+$ e $-$ da fonte.
  \item \textbf{Energizar} a placa com SW0 = 0, SW1 = 0 e SW9..7 = 000; depois a
        fonte (5\,V). Medir 3,3\,V no VCC do sensor e 5\,V no VCC do servo.
  \item \textbf{Programar:} \s{make jtag} (USB-Blaster na lista) e
        \s{make program}.
  \item \textbf{CP2102 no USB:} \s{ls /dev/ttyUSB*}; Processing ou
        \s{picocom}, um de cada vez.
\end{enumerate}
\end{minipage}\hfill
\begin{minipage}[t]{0.49\linewidth}
\vspace{0pt}
\textbf{Primeiro teste} (só avançar se a linha anterior passou)\\[2pt]
\begin{tabularx}{\linewidth}{@{}l >{\raggedright\arraybackslash}p{0.34\linewidth} Y@{}}
\toprule
& Como & Esperado \\
\midrule
EU1 & CP2102 sem FPGA, TXD ligado ao RXD; \s{picocom} & cada tecla volta na tela \\
--- & programado, SW1 = 0 & HEX5 = 0; LEDR4 e LEDR5 acesos (linhas em repouso) \\
EU2 & SW1 = 0, página 011; teclar \s{a}, \s{v}, \s{x} & HEX1 HEX0 = 61, 76, 78; paridade ok = 1; LEDR9 apagado \\
--- & SW1 = 1 & LEDR0 aceso; HEX5 HEX4 = \s{2 0}; servo em 020 \\
EU3 & Scope CH1 no \s{pwm} & 20\,ms; larguras da tabela do servo \\
EU4 & objeto a 10 e 30\,cm (régua), página 000 & distância $\pm$2\,cm; trigger de 10\,µs \\
EI1 & Processing & \s{aaa,ddd\#}; ângulos 020 $\to$ 160 $\to$ 020 \\
EI2 & teclar \s{a}, depois \s{v} & LEDR9 acende e o servo para no fim do ciclo; \s{v} retoma da posição mantida \\
\bottomrule
\end{tabularx}
\end{minipage}

\clearpage
% ==========================================================================
\section{Números-chave}

O ciclo dura cerca de 2\,s: espera, medida, envio de \s{aaa,ddd\#} e, no modo 0,
avanço do servo. Clock de 50\,MHz (20\,ns).

\begin{tabularx}{\linewidth}{@{}l l Y@{}}
\toprule
Grandeza & Valor & Origem \\
\midrule
Espera entre medidas & 2\,s = 100\,000\,000 clocks & \s{ESPERA} em \s{sonar2.v} \\
Timeout da medida & 10\,s = 500\,000\,000 clocks & \s{TIMEOUT} em \s{sonar2.v} \\
Pulso de trigger & 10\,µs = 500 clocks & \s{gerador\_pulso} \\
Echo $\to$ distância & 58,8\,µs por cm = 2941 clocks & \s{contador\_cm} (R = 2941) \\
Echo máximo real (4\,m) & cerca de 23,5\,ms & HC-SR04 \\
Sem eco (sensor desiste) & cerca de 47,5\,ms $\to$ lê 807\,cm & medido na bancada \bancada \\
Serial & 115200 bauds, 7 bits, paridade par, 1 stop (7E1) & TX e RX \\
Um caractere & 10 bits = 86,8\,µs & start + 7 + paridade + stop \\
Mensagem \s{aaa,ddd\#} & 8 caracteres, cerca de 700\,µs & ângulo, vírgula, distância, \s{\#} \\
PWM do servo & período 20\,ms (50\,Hz), pulso 0,7 a 2,2\,ms & \s{controle\_servo\_8} \\
Sincronizadores (\s{ligar}, \s{echo}, \s{entrada\_serial}) & 2 FFs cada, atraso de 40\,ns & \s{sincronizador.v} no \s{sonar2.v} \\
Simulação (testbenches) & espera 200\,µs, timeout 20\,ms & \s{sonar2\_tb\_config.vh} \\
\bottomrule
\end{tabularx}

Comandos: \s{a} = modo atenção (servo para), \s{v} = volta à varredura. Só
minúsculos e com paridade correta; o resto é ignorado. Desligado
(\s{ligar} = 0), os comandos são ignorados.

% ==========================================================================
\section{Chaves e LEDs}

\begin{minipage}[t]{0.40\linewidth}
\vspace{0pt}
\begin{tabularx}{\linewidth}{@{}l l Y@{}}
\toprule
Chave & Sinal & Uso \\
\midrule
SW0 & \s{reset} & 1 = reset assíncrono (estado 0, modo 0, posição 000) \\
SW1 & \s{ligar} & 1 = sonar funcionando; 0 = UC em \s{inicial} \\
SW9..7 & \s{sel\_depuracao} & página de HEX5..HEX0 \\
\bottomrule
\end{tabularx}
\end{minipage}\hfill
\begin{minipage}[t]{0.57\linewidth}
\vspace{0pt}
\begin{tabularx}{\linewidth}{@{}l l Y@{}}
\toprule
LED & Sinal & O que indica \\
\midrule
LEDR0 & \s{db\_ligar} & a chave SW1 direto (antes do sincronizador) \\
LEDR1 & \s{db\_echo} & echo do sensor \\
LEDR2 & \s{db\_trigger} & trigger (10\,µs, quase invisível) \\
LEDR3 & \s{db\_pwm} & PWM do servo (brilho fraco) \\
LEDR4 & \s{db\_saida\_serial} & linha TX; aceso em repouso \\
LEDR5 & \s{db\_entrada\_serial} & linha RX; aceso em repouso se o TXD está ligado \\
LEDR8..6 & \s{db\_posicao} & posição do servo em binário \\
LEDR9 & \s{db\_modo} & aceso fora da localização (atenção ou modo novo); o número do modo está no HEX4 da pág.\ 011 \\
\bottomrule
\end{tabularx}
\end{minipage}

% ==========================================================================
\section{Páginas dos displays (SW9..7)}

O multiplexador \s{U3\_PAGINA} (\s{sonar2.v}) escolhe 1 de 8 páginas de 6
dígitos. HEX5 é o display mais à esquerda.

\begin{tabularx}{\linewidth}{@{}l l Y Y Y Y Y@{}}
\toprule
SW9..7 & Página & HEX5 & HEX4 & HEX3 & HEX2 & HEX1 HEX0 \\
\midrule
000 & sonar (= exp5) & estado UC & estado HC-SR04 & posição 0--7 & distância (centena) & distância (dez., unid.) \\
001 & serial & estado TX & char TX (ASCII) & char TX (ASCII) & estado RX & char RX (ASCII) \\
010 & mensagem & ângulo & ângulo & ângulo & distância & distância \\
011 & comandos & estado UC & modo 0--3 (0 loc., 1 atenção, 2--3 livres) & índice do char 0--7 & paridade ok & char RX (ASCII) \\
1xx & livres (prova) & 0 & 0 & 0 & 0 & 0 0 \\
\bottomrule
\end{tabularx}

Simulação (caso 20): 000 = \s{201050}, 001 = \s{030078}, 010 = \s{040050},
011 = \s{200178}, 1xx = \s{000000}. ASCII: \s{a} = 61, \s{v} = 76, \s{A} = 41,
\s{x} = 78, \s{0}--\s{9} = 30--39, \s{,} = 2C, \s{\#} = 23.
Em funcionamento normal, a página 000 mostra \s{2 0} em HEX5 HEX4 quase o tempo
todo. Página nova: edite \s{s\_pagina4} a \s{s\_pagina7} em \s{sonar2.v}
(6 dígitos hex, 24 bits).

\clearpage
% ==========================================================================
\section{Placa: o que cada display, LED e chave mostra}

% Display de 7 segmentos: segmentos a..g na ordem do hexa7seg (1 = aceso)
\tikzset{
  segon/.style={line width=2pt, line cap=round, draw=red!80!black},
  segoff/.style={line width=2pt, line cap=round, draw=black!8},
}
\newcommand{\seg}[2]{\ifnum#1=1 \draw[segon] #2;\else\draw[segoff] #2;\fi}
\newcommand{\sete}[8]{%
  \begin{scope}[shift={#1}, scale=0.62, xslant=0.1]
    \draw[rounded corners=1.5pt, draw=black!35, fill=black!2] (-0.2,-0.2) rectangle (0.8,1.2);
    \seg{#2}{(0.12,1.0)--(0.48,1.0)}%
    \seg{#3}{(0.6,0.88)--(0.6,0.6)}%
    \seg{#4}{(0.6,0.4)--(0.6,0.12)}%
    \seg{#5}{(0.12,0.0)--(0.48,0.0)}%
    \seg{#6}{(0.0,0.4)--(0.0,0.12)}%
    \seg{#7}{(0.0,0.88)--(0.0,0.6)}%
    \seg{#8}{(0.12,0.5)--(0.48,0.5)}%
  \end{scope}}
\newcommand{\hexdig}[2]{\ifcase"#2 \sete{#1}1111110\or\sete{#1}0110000\or
  \sete{#1}1101101\or\sete{#1}1111001\or\sete{#1}0110011\or\sete{#1}1011011\or
  \sete{#1}1011111\or\sete{#1}1110000\or\sete{#1}1111111\or\sete{#1}1111011\or
  \sete{#1}1110111\or\sete{#1}0011111\or\sete{#1}1001110\or\sete{#1}0111101\or
  \sete{#1}1001111\or\sete{#1}1000111\fi}
% chave deslizante: #1 = posicao, #2 = 1 (para cima, ligada) ou 0
\newcommand{\chave}[2]{%
  \draw[rounded corners=1pt, draw=black!45, fill=black!4] ($#1+(-0.17,-0.32)$) rectangle ($#1+(0.17,0.32)$);
  \ifnum#2=1 \fill[black!70, rounded corners=0.8pt] ($#1+(-0.12,0.02)$) rectangle ($#1+(0.12,0.27)$);
  \else      \fill[black!70, rounded corners=0.8pt] ($#1+(-0.12,-0.27)$) rectangle ($#1+(0.12,-0.02)$);\fi}

Da esquerda para a direita, como na placa. Exemplo do desenho (caso 20 da simulação):
sonar ligado e em espera, servo na posição 001 (040°), última distância 050\,cm,
página 000. Os LEDs não dependem da página.

\begin{center}
\begin{tikzpicture}[
    ledon/.style={circle, draw=red!60!black, minimum size=3.4mm, inner sep=0pt, fill=red!75},
    ledoff/.style={circle, draw=red!60!black, minimum size=3.4mm, inner sep=0pt, fill=white},
    rot/.style={font=\scriptsize, align=center, text width=1.5cm, inner sep=1pt},
    tit/.style={font=\scriptsize\bfseries, text=black!65, inner sep=1pt}
  ]
  \def\L{1.62}
  % displays: HEX5 a HEX0
  \foreach \k [count=\i from 0] in {5,4,3,2,1,0}
    \node[tit] at (5.45+\i*0.8+0.2, -0.32) {HEX\k};
  \hexdig{(5.45+0*0.8,0)}{2}
  \hexdig{(5.45+1*0.8,0)}{0}
  \hexdig{(5.45+2*0.8,0)}{1}
  \hexdig{(5.45+3*0.8,0)}{0}
  \hexdig{(5.45+4*0.8,0)}{5}
  \hexdig{(5.45+5*0.8,0)}{0}
  \node[font=\scriptsize, text=black!65, anchor=west] at (10.4, 0.38) {página escolhida em SW9..7};
  % LEDs: LEDR9 a LEDR0
  \foreach \nome/\txt/\est [count=\i from 0] in {
      LEDR9/{modo $\neq$\\localização}/ledoff,
      LEDR8/{posição\\bit 2}/ledoff,
      LEDR7/{posição\\bit 1}/ledoff,
      LEDR6/{posição\\bit 0}/ledon,
      LEDR5/{entrada\\serial (RX)}/ledon,
      LEDR4/{saída\\serial (TX)}/ledon,
      LEDR3/{pwm}/ledoff,
      LEDR2/{trigger}/ledoff,
      LEDR1/{echo}/ledoff,
      LEDR0/{ligar\\(chave SW1)}/ledon} {
    \node[\est] at (0.4+\i*\L, -1.3) {};
    \node[tit] at (0.4+\i*\L, -0.95) {\nome};
    \node[rot, anchor=north] at (0.4+\i*\L, -1.55) {\txt};
  }
  % chaves: SW9 a SW0
  \foreach \nome/\txt/\pos [count=\i from 0] in {
      SW9/{página\\bit 2}/0, SW8/{página\\bit 1}/0, SW7/{página\\bit 0}/0,
      SW6/{livre}/0, SW5/{livre}/0, SW4/{livre}/0, SW3/{livre}/0, SW2/{livre}/0,
      SW1/{ligar}/1, SW0/{reset}/0} {
    \chave{(0.4+\i*\L, -2.95)}{\pos}
    \node[tit] at (0.4+\i*\L, -2.45) {\nome};
    \node[rot, anchor=north] at (0.4+\i*\L, -3.32) {\txt};
  }
\end{tikzpicture}
\end{center}

\textbf{O que cada display mostra em cada página} (mesmo exemplo; distância e ângulo em
decimal, caracteres em ASCII hexadecimal)

\begin{center}
\begin{tikzpicture}[
    pag/.style={draw=black!40, rounded corners=3pt, fill=blue!4, minimum width=2.75cm,
                minimum height=0.9cm, align=center, font=\footnotesize},
    lab/.style={font=\scriptsize, align=center, text width=#1, inner sep=1pt},
    lab/.default=1.45cm,
    br/.style={decorate, decoration={brace, mirror, amplitude=3pt}, draw=black!55},
    tit/.style={font=\scriptsize\bfseries, text=black!60, inner sep=1pt}
  ]
  \def\P{1.62}\def\X{3.55}
  % campo: #1 = y da linha, #2 = HEX da esquerda, #3 = HEX da direita, #4 = largura, #5 = texto
  \newcommand{\campo}[5]{\draw[br] ({\X+(5-#2)*\P-0.12}, #1-0.2) -- ({\X+(5-#3)*\P+0.5}, #1-0.2)
      node[midway, below=3pt, lab=#4] {#5};}
  \newcommand{\linha}[7]{% #1 = y; #2..#7 = digitos de HEX5 a HEX0
    \hexdig{(\X+0*\P,#1)}{#2}\hexdig{(\X+1*\P,#1)}{#3}\hexdig{(\X+2*\P,#1)}{#4}%
    \hexdig{(\X+3*\P,#1)}{#5}\hexdig{(\X+4*\P,#1)}{#6}\hexdig{(\X+5*\P,#1)}{#7}}
  \foreach \k [count=\i from 0] in {5,4,3,2,1,0}
    \node[tit] at (\X+\i*\P+0.2, 1.05) {HEX\k};

  \node[pag] at (1.4, 0.35) {\textbf{SW9..7 = 000}\\sonar (= exp5)};
  \linha{0}{2}{0}{1}{0}{5}{0}
  \campo{0}{5}{5}{1.5cm}{estado da UC\\(2 = espera)}
  \campo{0}{4}{4}{1.5cm}{estado do\\HC-SR04 (0)}
  \campo{0}{3}{3}{1.5cm}{posição 0--7\\(1 = 040°)}
  \campo{0}{2}{0}{4.2cm}{distância em cm (050)}

  \node[pag] at (1.4, -1.85) {\textbf{SW9..7 = 001}\\serial};
  \linha{-2.2}{0}{3}{0}{0}{7}{8}
  \campo{-2.2}{5}{5}{1.5cm}{estado do TX\\(0 = parado)}
  \campo{-2.2}{4}{3}{2.9cm}{caractere a enviar, ASCII\\(30 = '0')}
  \campo{-2.2}{2}{2}{1.5cm}{estado do RX\\(0 = parado)}
  \campo{-2.2}{1}{0}{2.9cm}{último caractere recebido,\\ASCII (78 = 'x')}

  \node[pag] at (1.4, -4.05) {\textbf{SW9..7 = 010}\\mensagem};
  \linha{-4.4}{0}{4}{0}{0}{5}{0}
  \campo{-4.4}{5}{3}{4.2cm}{ângulo em graus (040)}
  \campo{-4.4}{2}{0}{4.2cm}{distância em cm (050)}

  \node[pag] at (1.4, -6.25) {\textbf{SW9..7 = 011}\\comandos};
  \linha{-6.6}{2}{0}{0}{1}{7}{8}
  \campo{-6.6}{5}{5}{1.5cm}{estado da UC}
  \campo{-6.6}{4}{4}{1.5cm}{modo (0 loc.,\\1 atenção)}
  \campo{-6.6}{3}{3}{1.5cm}{índice do\\caractere 0--7}
  \campo{-6.6}{2}{2}{1.5cm}{paridade ok\\(1 = certa)}
  \campo{-6.6}{1}{0}{2.9cm}{último caractere recebido,\\ASCII (78 = 'x')}

  \node[pag] at (1.4, -8.45) {\textbf{SW9..7 = 1xx}\\livres};
  \linha{-8.8}{0}{0}{0}{0}{0}{0}
  \campo{-8.8}{5}{0}{9cm}{livres para a prova: editar \s{s\_pagina4} a \s{s\_pagina7} em \s{sonar2.v} (6 dígitos hexadecimais)}
\end{tikzpicture}
\end{center}

\begin{tabularx}{\linewidth}{@{}>{\raggedright\arraybackslash}p{0.27\linewidth} Y@{}}
\toprule
Código no display & Significado \\
\midrule
Estado da UC (HEX5, págs.\ 000 e 011) & 0 inicial · 1 preparação · 2 espera · 3 dispara medida · 4 espera medida · 9 timeout · 5 envia char · 6 espera tx · 7 próximo char · 8 muda posição · A final atenção · F final posição (B--E livres) \\
HC-SR04 (HEX4, pág.\ 000) & 0 inicial · 1 preparação · 2 envia trigger · 3 espera echo · 4 medida · 5 armazenamento · F final \\
TX (HEX5, pág.\ 001) · RX (HEX2, pág.\ 001) & TX: 0 inicial · 1 preparação · 3 espera · 7 transmissão · F final; RX: 0 inicial · 2 preparação · 3 espera · 7 recepção · 9 registra · F final \\
ASCII & \s{0}--\s{9} = 30--39 · \s{,} = 2C · \s{\#} = 23 · \s{a} = 61 · \s{v} = 76 · \s{x} = 78 · \s{A} = 41 · \s{V} = 56 \\
\bottomrule
\end{tabularx}

\clearpage
% ==========================================================================
\section{Códigos de estado}

\begin{center}
\begin{tikzpicture}[
    >={Stealth[length=2mm]},
    st/.style={draw=black!45, rounded corners=3pt, minimum width=2.45cm,
               minimum height=1.08cm, align=center, inner sep=1pt},
    bif/.style={st, draw=blue!60!black, very thick, fill=blue!7},
    lbl/.style={font=\scriptsize, text=black!65},
    modo/.style={font=\scriptsize\bfseries},
    ar/.style={->, draw=black!55}
  ]
  \newcommand{\est}[3]{{\bfseries\large #1}\\[-1pt]{\footnotesize #2}\\[-2pt]{\scriptsize\color{black!60}\ttfamily #3}}
  \def\dx{3.25}
  % linha de cima
  \node[st]  (s0) at (0*\dx, 0)   {\est{0}{inicial}{espera ligar}};
  \node[st]  (s1) at (1*\dx, 0)   {\est{1}{preparação}{zera}};
  \node[st]  (s2) at (2*\dx, 0)   {\est{2}{espera 2 s}{conta\_espera}};
  \node[st]  (s3) at (3*\dx, 0)   {\est{3}{dispara medida}{mede (1 clock)}};
  \node[st]  (s4) at (4*\dx, 0)   {\est{4}{espera medida}{conta\_timeout}};
  % linha do meio
  \node[st]  (s9) at (4*\dx, -2.1) {\est{9}{timeout}{limpa\_sensor}};
  \node[st]  (s5) at (3*\dx, -2.1) {\est{5}{envia char}{transmite}};
  \node[st]  (s6) at (2*\dx, -2.1) {\est{6}{espera tx}{fim\_transmissao}};
  \node[bif] (s7) at (1*\dx, -2.1) {\est{7}{próximo char}{conta\_char}};
  % fim do ciclo
  \node[st]  (s8) at (0*\dx, -4.0) {\est{8}{muda posição}{conta\_posicao}};
  \node[st]  (sA) at (1*\dx, -4.0) {\est{A}{final atenção}{fim\_posicao}};
  \node[st]  (sF) at (0*\dx, -5.7) {\est{F}{final posição}{fim\_posicao}};

  \coordinate (topy) at ($(s0.north)+(0,0.6)$);
  \coordinate (gx)   at ($(s0.west)+(-0.45,0)$);
  \coordinate (bot)  at ($(sF.south)+(0,-0.4)$);

  \draw[ar] (s0) -- (s1);
  \draw[ar] (s1) -- (s2);
  \draw[ar] (s2) -- (s3);
  \draw[ar] (s3) -- (s4);
  \node[lbl, above] at ($(s0.north east)!.5!(s1.north west)$) {ligar};
  \node[lbl, above] at ($(s2.north east)!.5!(s3.north west)$) {fim\_espera};

  \draw[ar] (s4.south) -- node[lbl, right] {10 s (fim\_timeout)} (s9.north);
  \draw[ar] ([xshift=-7mm]s4.south) -- ++(0,-0.45) -| node[lbl, pos=0.25, above] {fim\_medida} (s5.north);
  \draw[ar] (s9) -- (s5);
  \draw[ar] (s5) -- (s6);
  \draw[ar] (s6) -- (s7);
  \draw[ar] ([xshift=6mm]s7.south) -- ++(0,-0.45) -| node[lbl, pos=0.25, below] {faltam caracteres (8 no total)} (s5.south);

  \draw[ar] (s7.west) -| node[modo, pos=0.25, above] {modo 0} (s8.north);
  \draw[ar] (s7.south) -- node[modo, left] {modo 1} (sA.north);
  \draw[ar] (s8) -- (sF);

  \draw[ar] (sA.south) -- (sA.south |- bot) -- (gx |- bot) -- (gx |- topy)
            -- node[lbl, pos=0.5, above] {volta a 2: novo ciclo} (s2.north |- topy) -- (s2.north);
  \draw[draw=black!55] (sF.west) -- (gx |- sF.west);

  \node[align=left, anchor=north west, font=\footnotesize] at ($(s5.south west)+(-0.2,-1.3)$) {%
    \textbf{Regras da UC}\\
    \s{ligar} = 0: qualquer estado $\to$ 0 no clock seguinte\\
    modo (\s{Matual}, 2 bits): \s{a} $\to$ 1, \s{v} $\to$ 0, em qualquer estado\\
    códigos livres para modos novos: 2, 3\\
    em 0 (\s{inicial}) o modo volta a 0\\
    servo só anda em 8 (\s{conta\_posicao})\\
    códigos livres para estados novos: B, C, D, E};
\end{tikzpicture}
\end{center}

\textbf{Sinais de cada estado} (\s{sonar2\_uc}, Moore: cada sinal de controle
fica em 1 só no seu estado; nos demais, todos em 0).

\begin{tabularx}{\linewidth}{@{}l l >{\raggedright\arraybackslash}p{0.29\linewidth} >{\raggedright\arraybackslash}p{0.23\linewidth} Y@{}}
\toprule
Cód. & Estado & Sinal da UC em 1 $\to$ efeito no FD & Sai para (entrada da UC) & Nos pinos e LEDs \\
\midrule
0 & \s{inicial} & nenhum; modo forçado a 0 & 1 se \s{ligar} = 1 & servo parado na posição atual; \s{saida\_serial} = 1 (repouso) \\
1 & \s{preparacao} & \s{zera}: posição (U4) $\to$ 000, char (U6) $\to$ 0, espera (U8), TX e sensor reiniciados & 2 (1 clock) & \s{pwm} vai a 0,7\,ms (020°); LEDR8..6 = 000 \\
2 & \s{espera} & \s{conta\_espera}: U8 conta 2\,s & 3 com \s{fim\_espera} (U8) & tudo em repouso; servo se acomoda \\
3 & \s{dispara\_medida} & \s{mede}: U1 inicia a medida; zera o timeout (U9) & 4 (1 clock) & \s{trigger}: pulso de 10\,µs \\
4 & \s{espera\_medida} & \s{conta\_timeout}: U9 conta 10\,s & 5 com \s{fim\_medida} (pronto do U1); 9 com \s{fim\_timeout} (U9) & \s{echo} em 1 por 58,8\,µs/cm; HEX4 (sensor) 3 $\to$ 4 $\to$ 5 \\
9 & \s{timeout\_medida} & \s{limpa\_sensor}: reinicia U1 (sai de espera echo) & 5 (1 clock) & mensagem sai com distância 000 \\
5 & \s{envia\_char} & \s{transmite}: partida do TX (U2) com o char do mux (U7) & 6 (1 clock) & \s{saida\_serial} desce (start bit) \\
6 & \s{espera\_tx} & nenhum & 7 com \s{fim\_transmissao} (pronto do U2) & \s{saida\_serial}: quadro de 86,8\,µs \\
7 & \s{proximo\_char} & \s{conta\_char}: U6 avança o índice (0--7) & 5 se falta char; com \s{fim\_char}: 8 (modo 0) ou A (modo 1) & bifurcação; HEX3 da pág.\ 011 = índice \\
8 & \s{muda\_posicao} & \s{conta\_posicao}: U4 avança (111 $\to$ 000) & F (1 clock) & \s{pwm} muda de largura: servo anda 20°; LEDR8..6 \\
F & \s{final\_posicao} & \s{fim\_posicao} & 2 (1 clock) & pino \s{fim\_posicao}: pulso de 20\,ns \\
A & \s{final\_atencao} & \s{fim\_posicao} & 2 (1 clock) & pulso de 20\,ns; servo não anda \\
\bottomrule
\end{tabularx}

Sempre ativos, fora da UC: \s{pwm} (U3, 20\,ms, largura pela posição), recepção
(U10) e FSM de modo. \s{ligar} = 0 leva a 0 de qualquer estado; \s{reset} leva a 0
e ao modo 0. Saídas do topo: \s{trigger}, \s{pwm}, \s{saida\_serial},
\s{fim\_posicao}, \s{db\_modo} (LEDR9), HEX5..0 e LEDR8..0.

\begin{tabularx}{\linewidth}{@{}l l Y@{}}
\toprule
Máquina & Onde aparece & Códigos \\
\midrule
HC-SR04 (\s{interface\_hcsr04\_uc}) & HEX4, pág.\ 000 & 0 inicial · 1 preparação · 2 envia trigger · 3 espera echo · 4 medida · 5 armazenamento · F final \\
TX (\s{tx\_serial\_uc}) & HEX5, pág.\ 001 & 0 inicial · 1 preparação · 3 espera · 7 transmissão · F final \\
RX (\s{rx\_serial\_uc}) & HEX2, pág.\ 001 & 0 inicial · 2 preparação · 3 espera · 7 recepção · 9 registra · F final \\
\bottomrule
\end{tabularx}

% ==========================================================================
\section{Servo: posição, ângulo e PWM}

\begin{minipage}[t]{0.62\linewidth}
\vspace{0pt}
\begin{tabularx}{\linewidth}{@{}l l l l Y@{}}
\toprule
Posição & Ângulo & Largura & Clocks & ASCII da ROM \\
\midrule
000 & 020 & 0,700\,ms & 35\,000  & 30 32 30 \\
001 & 040 & 0,914\,ms & 45\,700  & 30 34 30 \\
010 & 060 & 1,129\,ms & 56\,450  & 30 36 30 \\
011 & 080 & 1,343\,ms & 67\,150  & 30 38 30 \\
100 & 100 & 1,557\,ms & 77\,850  & 31 30 30 \\
101 & 120 & 1,771\,ms & 88\,550  & 31 32 30 \\
110 & 140 & 1,986\,ms & 99\,300  & 31 34 30 \\
111 & 160 & 2,200\,ms & 110\,000 & 31 36 30 \\
\bottomrule
\end{tabularx}
\end{minipage}\hfill
\begin{minipage}[t]{0.34\linewidth}
\vspace{0pt}
A posição (contador U4) escolhe a largura do pulso e o ângulo da ROM. LEDR8..6
mostram a posição em binário.

No modo 0 a varredura vai 000 $\to$ 111 e volta a 000 (dente de serra).

Período de 20\,ms (1\,000\,000 clocks). Larguras em
\s{controle\_servo\_8.v}, ângulos em \s{rom\_angulos\_8x24.v}.
\end{minipage}

% ==========================================================================
\section{Analog Discovery / WaveForms}

Canais conforme o planejamento. Todos os sinais em 3,3\,V do lado da FPGA; GND
do AD no GND da placa.

\begin{tabularx}{\linewidth}{@{}l l Y Y@{}}
\toprule
Sinal & Canal & Ferramenta e configuração & O que esperar \\
\midrule
\s{entrada\_serial} & DIO6 & Logic + UART ou Protocol: 115200, 7 bits, Even, Standard, LSB first; trigger $\downarrow$; 20\,µs/div & quadro de 86,8\,µs; \s{a} = 61, \s{v} = 76 \\
\s{saida\_serial} & DIO7 & mesma configuração; 100\,µs/div & \s{aaa,ddd\#} em cerca de 700\,µs \\
\s{pwm} & CH1+ & Scope: subida, 1,5\,V; 500\,µs/div (largura) ou 5\,ms/div (período) & 20\,ms; larguras da tabela do servo \\
\s{trigger} & CH2+ & Scope: captura única, subida; 5\,µs/div & pulso de 10\,µs \\
\s{echo} & DIO0 & Logic: subida; 2\,ms/div (10\,ms/div para o caso sem eco) & 58,8\,µs/cm; cerca de 47\,ms quando lê 807 \\
\s{fim\_posicao} & DIO1 & Logic: 100\,MHz; modo Record para vários ciclos & pulso de 20\,ns a cada 2\,s \\
\bottomrule
\end{tabularx}

\textbf{Trigger e echo juntos no Scope} (usado na bancada em 08/10): mostra se o
sensor responde ao disparo e a largura do echo. Mudar o 1$+$ do \s{pwm} para o
ECHO (JP2 · 4); o 2$+$ continua no TRIG.

\begin{tabularx}{\linewidth}{@{}l Y l Y@{}}
\toprule
Campo & Valor & Campo & Valor \\
\midrule
Mode / Trigger & Repeated / Auto & Time Base & 5,8\,ms/div \\
Source & Channel 1 (\s{echo}) & Time Position & 4,971\,ms (disparo perto da esquerda) \\
Type / Condition & Edge / Rising & Channel 1 e 2 & Offset 0\,V, Range 1\,V/div \\
Level / Hyst. / HoldOff & 2\,V / Auto / 0\,s & Amostragem & 8192 amostras a 141\,kHz \\
\bottomrule
\end{tabularx}

O que se vê: o trigger (CH2, azul) é um traço fino logo antes da subida do echo
(10\,µs some a 5,8\,ms/div; para medir, 5\,µs/div); o echo (CH1, laranja) é um
pulso de cerca de 3,3\,V com 58,8\,µs por cm (3,4\,ms $\approx$ 57\,cm). A janela
vai até cerca de 34\,ms depois do disparo: cobre 4\,m (23,5\,ms), mas não o caso sem
eco (47\,ms; usar 10\,ms/div). Echo em 0 depois do traço do trigger = sensor não
respondeu (alimentação, GND ou fio do ECHO).

\begin{itemize}
  \item \textbf{Texto decodificado:} no UART do Logic, Format = ASCII.
        \textbf{Congelar:} Mode = Single. \textbf{Salvar figura:} File $\to$ Export.
  \item \textbf{Comando sem o CP2102:} Protocol $\to$ UART, TX no DIO6, Send \s{a}
        (antes, desligar o TXD do CP2102 do JP1 pino 4).
  \item \textbf{Servo no tempo:} Data Logger, CH1 no \s{pwm}, função \s{C1DC / 3.3 * 20}
        (largura em ms): escada da varredura e patamar no modo atenção.
  \item \textbf{Distância:} Scope com o \s{echo} num canal analógico, Measurements
        $\to$ PosWidth; Custom \s{Channel.measure("PosWidth", T1, T2) / 58e-6} dá cm.
  \item Scope e Data Logger usam os mesmos canais analógicos: um de cada vez.
\end{itemize}

% ==========================================================================
\section{Sintoma $\to$ causa $\to$ ação}

Ordem: alimentação e fios, chaves e páginas, Analog Discovery, comparação com a
simulação. \bancada{} = aconteceu na bancada.

\begin{tabularx}{\linewidth}{@{}>{\raggedright\arraybackslash}p{0.30\linewidth} Y Y@{}}
\toprule
Sintoma & Causa provável & Ação \\
\midrule
Nada acontece, HEX5 = 0 & SW1 = 0 ou SW0 = 1 & Ajustar as chaves; LEDR0 aceso = ligar ok \\
Mensagens chegam, mas \s{a} não muda nada \bancada & Fio TXD $\to$ JP1 pino 4 solto ou no pino errado & LEDR5 aceso em repouso; página 011 deve mostrar 61; trocar o jumper \\
\s{a} chega (61), paridade ok = 0 & Terminal fora de 7E1 & 115200, 7 bits, paridade par \\
LEDR9 acende, mas o servo anda mais uma posição & \s{a} chegou depois da bifurcação & Previsto: para no fim do ciclo seguinte \\
Comando ignorado com o sonar desligado & Previsto: em \s{inicial} o modo fica em 0 & Ligar antes de enviar \\
HEX5 HEX4 = \s{3 0} parado após desligar/religar SW1 \bancada & SW1 assíncrona + UC one-hot sem safe state $\to$ estado ilegal & Corrigido com 2 FFs em \s{ligar}, \s{echo} e \s{entrada\_serial} (\s{sincronizador.v}); se voltar, reset (SW0) \\
Leitura de 807\,cm em algumas posições \bancada & Sensor sem eco: segura o echo cerca de 47\,ms & Objeto plano e rígido a menos de 4\,m; Processing mostra ``Não Detectado'' \\
HEX5 = 4, HEX4 = 3 por 10\,s, depois \s{aaa,000\#} & Sensor não respondeu (timeout) & Conferir 3,3\,V (JP2 · 29), GND e fio do echo \\
HEX5 = 6 preso & Falha na transmissão & Analisar \s{saida\_serial} no DIO7 \\
Servo parado ou tremendo & Queda de tensão do servo & Verificar 5\,V da fonte externa e o $-$ dela no GND comum \\
\s{Device or resource busy} & Outro programa usa a porta & Fechar Processing ou picocom \\
Processing não envia a tecla & Foco no editor & Clicar na janela do radar \\
\s{A} maiúsculo não funciona & Caps Lock: \s{A} = 41 é ignorado & Desligar Caps Lock \\
\s{No JTAG hardware available} & Cabo USB da placa desligado & Ligar o cabo; \s{make jtag} \\
\bottomrule
\end{tabularx}

% ==========================================================================
\section{Comandos}

\begin{minipage}[t]{0.50\linewidth}
\vspace{0pt}
\begin{tabularx}{\linewidth}{@{}l Y@{}}
\toprule
Comando (em \s{labdig2/exp6}) & O que faz \\
\midrule
\s{make check TB=sonar2\_tb} & Icarus (cerca de 1\,min) \\
\s{make simb TB=sonar2\_tb} & Questa batch; ``todos os casos OK'' \\
\s{make sim TB=sonar2\_tb} & Questa com ondas \\
\s{sim cenario\_m\_tb} & igual a \s{make sim TB=...}; aceita o módulo ou \s{tb/cenario\_m\_tb.v} (Tab), não \s{cenario\_m\_tb.v}; \s{simb}/\s{check} direto têm limite de 60\,s (\s{sonar2\_tb} no Icarus: use \s{make check}) \\
\s{make regressao} & todos os \s{*\_tb.v} no Questa, com resumo \\
\s{make} & compila (\s{.sof} + timing) \\
\s{make timing} & relatório de timing \\
\s{make jtag} / \s{make program} & lista o USB-Blaster / grava \\
\s{make gui} & Quartus (Pin Planner, RTL Viewer) \\
\s{make entrega} & \s{.qar} e zip do ModelSim em \s{entrega/} \\
\s{make entrega-teste} & roda o zip extraído (antes de entregar) \\
\bottomrule
\end{tabularx}
\end{minipage}\hfill
\begin{minipage}[t]{0.47\linewidth}
\vspace{0pt}
\textbf{Terminal serial} (fechar o Processing antes):\\
\s{ls /dev/ttyUSB*}\\
\s{picocom -b 115200 -d 7 -y e -p 1 /dev/ttyUSB0}

\s{-y} = paridade, \s{-p} = stop bits. Ctrl+A Ctrl+X sai; Ctrl+A Ctrl+C liga o
eco local. Teste de paridade errada: 8N1 (\s{-d 8 -y n}) e teclar \s{v}: o modo
não muda.

\textbf{Processing:} porta \s{/dev/ttyUSB0}, já em 7E1. Clicar na janela do
radar e teclar; o console mostra \s{Enviando tecla 'a' ...}. Acima de 100\,cm
aparece ``Não Detectado''; só abaixo de 50\,cm vira ponto no radar.
\end{minipage}

\clearpage
% ==========================================================================
\section{Prova (modelo P1): roteiro por atividade}

A prova não é um ciclo de desenvolvimento: são três atividades, cada uma
terminando numa apresentação ao professor. Os workflows das duas páginas
seguintes entram dentro delas (síntese e placa na 1, simulação na 2). A
Atividade 3 é projeto no papel. Cada seta: apresentou ao professor; esperando por ele, adiante a próxima.

\begin{center}
\begin{tikzpicture}[
    >={Stealth[length=2mm]},
    st/.style={draw=black!45, rounded corners=3pt, minimum width=3.15cm,
               minimum height=1.05cm, align=center, inner sep=1pt},
    fim/.style={st, draw=green!40!black, thick, fill=green!6},
    lbl/.style={font=\scriptsize, text=black!70, inner sep=1pt},
    ar/.style={->, draw=black!60}
  ]
  \newcommand{\wf}[3]{{\bfseries #1}\,·\,#2\\[-1pt]{\scriptsize\color{black!65}#3}}
  \def\dx{3.75}
  \node[st]  (P0) at (0*\dx, 0) {\wf{0}{bancada}{montar + sintetizar}};
  \node[st]  (P1) at (1*\dx, 0) {\wf{1}{Atividade 1}{placa, testes, demo}};
  \node[st]  (P2) at (2*\dx, 0) {\wf{2}{Atividade 2}{cenário M, ondas}};
  \node[st]  (P3) at (3*\dx, 0) {\wf{3}{Atividade 3}{projeto no papel}};
  \node[fim] (PF) at (4*\dx, 0) {\wf{F}{relato}{até o fim do dia}};
  \draw[ar] (P0) -- (P1);
  \draw[ar] (P1) -- (P2);
  \draw[ar] (P2) -- (P3);
  \draw[ar] (P3) -- (PF);
  \draw[ar] (P1.south) to[bend right=18] node[lbl, below] {esperando o professor: adiantar a 2 e a 3} (P3.south);
\end{tikzpicture}
\end{center}

\begin{tabularx}{\linewidth}{@{}>{\raggedright\arraybackslash}p{0.14\linewidth} Y@{}}
\toprule
Item & O que fazer e mostrar \\
\midrule
0 bancada & montagem da pág.\ 1 (GND primeiro; sensor em 3,3\,V; servo na fonte de 5\,V); enquanto isso, síntese: \s{make} (PC: Ctrl+L) \\
1a, b, c & \s{make program}; plano de testes do Planejamento: EU1--EU4, EI1--EI2 e EI3 = ET1--ET12; demonstração: varredura, \s{a} (LEDR9, servo para), \s{v} (retoma), páginas \\
1e-1: caracteres $\to$ modo & \s{sonar2.v:129} sincronizador da \s{entrada\_serial} $\to$ \s{sonar2\_fd.v:235} receptor U10 (pronto, dado, paridade) $\to$ \s{sonar2\_fd.v:251} comparadores: \s{recebeu\_atencao = pronto \& paridade\_ok \& (dado == 'a')} $\to$ \s{sonar2\_uc.v:101} FSM de modo (\s{Matual}) $\to$ \s{sonar2\_uc.v:139} bifurcação no fim do ciclo. Outro caractere: a UC nem fica sabendo \\
1e-2: estado da transmissão & página 001 (\s{sonar2.v:224}): HEX5 = estado do TX (0 inicial, 1 preparação, 3 espera, 7 transmissão, F final; \s{tx\_serial\_uc.v:32}), HEX4 HEX3 = caractere em ASCII; pág.\ 011 HEX3 = índice 0--7; LEDR4 = \s{saida\_serial}; AD: DIO7 com Protocol UART 7E1; simulação: grupo Mensagem\_TX do \s{wave.do} \\
2f, g & o Cenário M já está em \s{tb/cenario\_m\_tb.v} (020, 040, 060; \s{a} na espera; 2 ciclos em atenção, em 080). Para o servo parar em 060: \s{espera\_trigger; envia\_comando("a", ATENCAO); ciclo(2, 35, PARADO);} \\
2 rodar & \s{simb tb/cenario\_m\_tb.v} (confere) e \s{sim tb/cenario\_m\_tb.v} (ondas). PC: \s{do simula.do cenario\_m\_tb} \\
2h & colar o \s{wave zoom range} do relatório ``Troca de modo 1''. Grupo Troca\_de\_modo: \s{rx\_pronto}, \s{rx\_dado} = a, \s{paridade\_ok}, \s{recebeu\_atencao} (condição do FD) e \s{modo} (\s{Matual}) mudando; \s{estado}: 7 \s{proximo\_char} $\to$ 8 \s{muda\_posicao} $\to$ F \s{final\_posicao} $\to$ 2 \s{espera} $\to$ 3 \s{dispara\_medida}. Print para o relato \\
3k, l (receita) & (1) interpretar: comando que entra, que sai, comportamento, quando vale (fim do ciclo). (2) FD: comparadores dos caracteres novos e componentes novos (contador) só com controle vindo da UC. (3) UC: FSM de modo + ramo na bifurcação + estados novos (B--E) + tabela de saídas. (4) anotar os diagramas impressos e fotografar. (5) mostrar que a restrição vale: nada reseta nem inibe componente pelo FD \\
3 modelo (\s{d}/\s{x}) & FD: \s{recebeu\_duplo} (\s{d}) e \s{recebeu\_normal} (\s{x}), iguais aos de \s{a}/\s{v}; contador U11 de 1\,s (\s{conta\_espera\_duplo} $\to$ \s{fim\_espera\_duplo}). Modo: LOCALIZACAO $\to$ DUPLO (código 2) com \s{d}; DUPLO $\to$ LOCALIZACAO com \s{x}. Bifurcação em \s{proximo\_char}: DUPLO e 1ª medida $\to$ \s{espera\_duplo} (B, 1\,s) $\to$ \s{dispara\_medida}; DUPLO e 2ª $\to$ \s{muda\_posicao}. Flag \s{segunda} na UC: 1 em \s{espera\_duplo}, 0 em \s{muda\_posicao} e \s{inicial}. Validado em simulação (08/10): \s{prova/modelo\_duplo.patch} \\
F relato & foto dos diagramas anotados, print das ondas do item h, resultado do plano de testes; enviar até o fim do dia \\
\bottomrule
\end{tabularx}

\clearpage
% ==========================================================================
\section{Testbench na prova: do enunciado às ondas (Atividade 2)}

\begin{center}
\begin{tikzpicture}[
    >={Stealth[length=2mm]},
    st/.style={draw=black!45, rounded corners=3pt, minimum width=2.15cm,
               minimum height=1.0cm, align=center, inner sep=1pt},
    err/.style={st, draw=red!60!black, thick, fill=red!6},
    fim/.style={st, draw=green!40!black, thick, fill=green!6},
    lbl/.style={font=\scriptsize, text=black!70, inner sep=1pt},
    lble/.style={lbl, text=red!60!black},
    ar/.style={->, draw=black!60},
    are/.style={->, draw=red!60!black}
  ]
  \newcommand{\wf}[3]{{\bfseries #1}\,·\,#2\\[-1pt]{\scriptsize\color{black!65}#3}}
  \def\dx{2.6}
  \node[st]  (T0) at (0*\dx, 0) {\wf{0}{lê}{o que, quando}};
  \node[st]  (T1) at (1*\dx, 0) {\wf{1}{traduz}{frase $\to$ linha}};
  \node[st]  (T2) at (2*\dx, 0) {\wf{2}{escreve}{cp + módulo}};
  \node[st]  (T3) at (3*\dx, 0) {\wf{3}{roda}{\ttfamily simb <tb>}};
  \node[st]  (T5) at (4*\dx, 0) {\wf{5}{ondas}{\ttfamily sim <tb>}};
  \node[st]  (T6) at (5*\dx, 0) {\wf{6}{localiza}{zoom range}};
  \node[fim] (TF) at (6*\dx, 0) {\wf{F}{apresenta}{log + ondas}};
  \node[err] (T4) at (3*\dx, -1.55) {\wf{4}{conserta}{tabela abaixo}};
  \draw[ar] (T0) -- (T1);
  \draw[ar] (T1) -- (T2);
  \draw[ar] (T2) -- (T3);
  \draw[ar] (T3) -- (T5);
  \draw[ar] (T5) -- (T6);
  \draw[ar] (T6) -- (TF);
  \draw[are] ([xshift=-4mm]T3.south) -- node[lble, left] {ERRO, FALHA} ([xshift=-4mm]T4.north);
  \draw[ar] ([xshift=4mm]T4.north) -- node[lbl, right] {corrigido} ([xshift=4mm]T3.south);
\end{tikzpicture}
\end{center}

\textbf{1 traduz: cada frase do cenário vira uma ou duas linhas} (comece cada uma com
\s{novo\_caso(n, "frase")}: o log e as ondas ficam com o texto do enunciado)

\begin{tabularx}{\linewidth}{@{}>{\raggedright\arraybackslash}p{0.30\linewidth} Y@{}}
\toprule
Frase do enunciado & Linhas \\
\midrule
liga o circuito & \s{inicia;} \\
$N$ ciclos em localização (posições 20° a \ldots) & \s{ciclo(0, cm, ANDA); ciclo(1, cm, ANDA); \ldots} (posição $p$ = (ângulo $-$ 20)/20) \\
recebe \s{a} (sem dizer quando) & na espera: \s{envia\_comando("a", ATENCAO);} e depois o \s{ciclo} \\
\s{a} durante a medida & \s{espera\_trigger; envia\_comando("a", ATENCAO); ciclo(p, cm, PARADO);} \\
\s{a} durante a transmissão & \s{responde\_echo(cm); espera\_transmissao; envia\_comando("a", ATENCAO); fim\_ciclo(p, cm, PARADO);} \\
$M$ ciclos em atenção ($M$ medidas) & \s{ciclo(p, cm, PARADO);} $M$ vezes, mesma posição \\
recebe \s{v} & \s{envia\_comando("v", LOCALIZACAO);} e o próximo \s{ciclo(p, cm, ANDA)} \\
caractere inválido, maiúsculo & \s{envia\_comando("A", modo\_atual);} (o modo não muda) \\
paridade errada & \s{envia\_comando\_paridade\_errada("v", modo\_atual);} \\
sensor não responde (sem objeto) & \s{ciclo\_sem\_echo(p, ANDA ou PARADO);} \\
desliga e religa · reset & \s{desliga; confere\_modo(LOCALIZACAO); liga;} · \s{aplica\_reset;} (depois: posição 0) \\
duas medidas na mesma posição & \s{medida(p, cm1); medida(p, cm2); espera\_fim\_ciclo; confere\_posicao(...);} \\
intervalo de tempo & \s{t0 = \$time;} \ldots\ \s{confere\_tempo("nome", t\_trigger - t0, min, max);} \\
mostrar no display ou LED & \s{confere\_pagina(3'b011, 24'h...);} · \s{confere\_sinal("LEDR9", db\_modo\_out, 1);} \\
\bottomrule
\end{tabularx}

\textbf{Para prever o esperado:} ângulo = 20 + 20$p$; mensagem \s{aaa,ddd\#} com
\s{ddd} = cm (de 1 a 300). Em localização o servo vai para $p+1$ depois do ciclo (7 $\to$
0); em atenção fica. Comando que chega antes do fim do ciclo (espera, medida ou
transmissão) já vale no próprio ciclo. Depois de \s{inicia}, religar ou reset: modo
localização e posição 0. Um ciclo leva cerca de 1 a 2\,ms; o \s{envia\_comando}, 87\,µs.

\begin{tabularx}{\linewidth}{@{}>{\raggedright\arraybackslash}p{0.33\linewidth} Y Y@{}}
\toprule
4 conserta: mensagem & Causa provável & O que fazer \\
\midrule
\s{ERRO: servo na posicao X (esperada Y)} & \s{ANDA}/\s{PARADO} trocado, ou o comando chegou em outro ciclo & regra do fim do ciclo; \s{rastro\_uc = 1;} \\
\s{ERRO: mensagem "..." esperada "..."} & posição ou cm errados no \s{ciclo} & recalcular: ângulo = 20 + 20$p$ \\
\s{ERRO: UC no modo X (esperado Y)} & caractere errado, minúsculo/maiúsculo, comando que não vale naquele modo & conferir o caractere e o modo anterior \\
\s{FALHA: simulacao travada (watchdog)} & esperando algo que não acontece (estado que não vem, ciclo com o sonar desligado) & \s{rastro\_uc = 1;} e ver onde parou \\
\s{Unable to bind} · \s{Undefined variable} & nome errado ou variável declarada depois de usar & declarar no começo do módulo \\
nada compila & erro de sintaxe neste ou em outro arquivo de \s{tb/}; dois arquivos com o mesmo \s{module} & corrigir; \s{do simula.do <tb>} isola \\
\bottomrule
\end{tabularx}

\textbf{5--6 nas ondas} (\s{sim tb/cenario\_x\_tb.v}; PC: \s{do simula.do cenario\_x\_tb}):
o \s{wave.do} abre os grupos Testbench (\s{caso}, \s{caso\_texto}),
Troca\_de\_modo (\s{entrada\_serial}, \s{rx\_pronto}, \s{rx\_dado}, \s{paridade\_ok},
\s{recebeu\_*}, \s{modo}, \s{estado}), entradas e saídas, condições e controles da UC,
mensagem TX, sensor e servo; estados e modos pelo nome.
\begin{itemize}[before=\raggedright]
  \item \textbf{Troca de modo:} o \s{finaliza} imprime ``== Troca de modo'' com os sinais
        do FD, a troca do \s{Matual} e os estados da UC; copiar a linha
        \s{wave zoom range ...} e colar no console do Questa.
  \item \textbf{Outro evento} (timeout, comando ignorado, mudança de posição):
        \s{t0 = \$time;} antes e \s{relata\_janela("nome", t0, \$time);} depois (mesmo
        relatório, para qualquer janela); \s{marca("texto");} registra um instante.
  \item \textbf{Mostrar:} cursores (Insert Cursor) na subida de \s{recebeu\_*} e na troca
        de \s{modo}: o Questa mostra o tempo entre eles. Figura: File $\to$ Export $\to$ Image.
  \item \textbf{Apresentar:} o cenário comentado (um \s{novo\_caso} por frase), o log com
        ``todos os casos OK'' e a janela das ondas com os sinais pedidos.
\end{itemize}

\clearpage
% ==========================================================================
\section{Workflow como máquina de estados: notebook (Linux)}

Máquina de Moore: em cada estado você executa a saída (os comandos da tabela) e
o resultado escolhe a transição. Seta preta: deu certo. Seta vermelha: deu
errado, vai para \textbf{5 corrige}, que devolve ao circuito ou ao testbench.
Terminal sempre na pasta \s{exp6}; \s{<tb>} = módulo ou caminho do arquivo.

\begin{center}
\begin{tikzpicture}[
    >={Stealth[length=2mm]},
    st/.style={draw=black!45, rounded corners=3pt, minimum width=2.6cm,
               minimum height=0.95cm, align=center, inner sep=1pt},
    err/.style={st, draw=red!60!black, thick, fill=red!6},
    fim/.style={st, draw=green!40!black, thick, fill=green!6},
    lbl/.style={font=\scriptsize, text=black!70, inner sep=1pt},
    lble/.style={lbl, text=red!60!black},
    ar/.style={->, draw=black!60},
    are/.style={->, draw=red!60!black}
  ]
  \newcommand{\wf}[3]{{\bfseries #1}\,·\,#2\\[-1pt]{\scriptsize\color{black!65}\ttfamily #3}}
  \def\dx{3.4}\def\dy{1.75}
  \node[st]  (L0) at (0*\dx, 0)      {\wf{0}{início}{pasta exp6}};
  \node[st]  (L1) at (1*\dx, 0)      {\wf{1}{planeja}{grep EXTENSAO}};
  \node[st]  (L2) at (2*\dx, 0)      {\wf{2}{circuito}{rtl/ · make sync}};
  \node[st]  (L3) at (3*\dx, 0)      {\wf{3}{testbench}{config · cenário}};
  \node[st]  (L4) at (4*\dx, 0)      {\wf{4}{simula}{simb <tb>}};
  \node[err] (L5) at (3*\dx, -\dy)   {\wf{5}{corrige}{sim <tb> · rs}};
  \node[st]  (L6) at (4*\dx, -2*\dy) {\wf{6}{regressão}{make regressao}};
  \node[st]  (L7) at (3*\dx, -2*\dy) {\wf{7}{síntese}{make · make timing}};
  \node[st]  (L8) at (2*\dx, -2*\dy) {\wf{8}{placa}{make program}};
  \node[st]  (L9) at (1*\dx, -2*\dy) {\wf{9}{entrega}{make entrega-teste}};
  \node[fim] (LF) at (0*\dx, -2*\dy) {\wf{F}{fim}{.qar e .zip}};

  \draw[ar] (L0) -- (L1);
  \draw[ar] (L1) -- (L2);
  \draw[ar] (L2) -- (L3);
  \draw[ar] (L3) -- (L4);
  \draw[ar] (L1.north) to[bend left=22] node[lbl, above] {só cenário} (L3.north);
  \draw[ar] (L4) -- (L6);
  \draw[ar] (L6) -- (L7);
  \draw[ar] (L7) -- (L8);
  \draw[ar] (L8) -- (L9);
  \draw[ar] (L9) -- (LF);
  \draw[are] (L4.south west) -- node[lble, pos=0.5, right=3pt] {FALHA, erro} (L5.north east);
  \draw[are] (L6.north west) -- node[lble, pos=0.5, right=3pt] {FALHA} (L5.south east);
  \draw[are] (L7) -- node[lble, right] {erro} (L5);
  \draw[are] (L8.north east) -- node[lble, pos=0.5, left=3pt] {placa $\neq$ simulação} (L5.south west);
  \draw[ar] ([xshift=3mm]L5.north) -- node[lbl, right] {cenário} ([xshift=3mm]L3.south);
  \draw[ar] (L5.north west) -- node[lbl, pos=0.5, left=3pt] {circuito} (L2.south east);
  \draw[ar] (LF.north) -- ++(0,0.6) -| node[lbl, pos=0.25, above] {próximo pedido} (L1.south);
\end{tikzpicture}
\end{center}

\begin{tabularx}{\linewidth}{@{}>{\raggedright\arraybackslash}p{0.13\linewidth} Y >{\raggedright\arraybackslash}p{0.22\linewidth}@{}}
\toprule
Estado & O que fazer & Vai para \\
\midrule
0 início & terminal na pasta: \s{cd \textasciitilde/projetos/labdig2/exp6} & 1 \\
1 planeja & ler o pedido; \s{grep -n EXTENSAO rtl/*.v *.vh wave.do}; guia de modificações da cola & 2 se muda o circuito; 3 se é só cenário \\
2 circuito & editar \s{rtl/} nos pontos \s{EXTENSAO}; arquivo novo em \s{rtl/}: \s{make sync} & 3 \\
3 testbench & circuito mudou: \s{sonar2\_tb\_config.vh}; cenário: \s{cp tb/cenario\_m\_tb.v tb/cenario\_x\_tb.v} e trocar o nome do módulo & 4 \\
4 simula & \s{simb <tb>} (ou \s{make simb TB=cenario\_x\_tb}); 6\,s por cenário & 6 se ``todos os casos OK''; 5 se FALHA ou erro de compilação \\
5 corrige & \s{sim <tb>}: ondas com o \s{wave.do}; colar o \s{wave zoom range} do relatório; \s{rastro\_uc = 1;} no cenário; editar e digitar \s{rs} no Questa & 2 se o erro é do circuito; 3 se é do cenário \\
6 regressão & \s{make regressao}: todos os \s{*\_tb.v}, resumo no fim & 7 se 0 com falha; 5 se não \\
7 síntese & \s{make} e \s{make timing} & 8 se TIMING PASS; 5 se erro \\
8 placa & \s{make jtag}, \s{make program}; SW0 = reset, SW1 = ligar, SW9..7 = página; tabela sintoma $\to$ causa $\to$ ação & 9 se funciona; 5 se a placa difere da simulação \\
9 entrega & \s{make entrega-teste} (zip extraído roda?); depois \s{make entrega} (\s{entrega/}) & F se OK; 5 se o zip falha \\
F fim & enviar o \s{.qar} e o \s{.zip}; para o PC da faculdade: \s{make pacote} & 1 no próximo pedido \\
\bottomrule
\end{tabularx}

\clearpage
% ==========================================================================
\section{Workflow como máquina de estados: PC da faculdade}

Sem \s{make}: Quartus e ModelSim/Questa pela interface. A pasta vem do notebook
(\s{make pacote}): o Quartus abre o projeto dela e o ModelSim simula a partir
dela, então os dois usam a mesma cópia do RTL. Mesma convenção de setas.

\begin{center}
\begin{tikzpicture}[
    >={Stealth[length=2mm]},
    st/.style={draw=black!45, rounded corners=3pt, minimum width=2.6cm,
               minimum height=0.95cm, align=center, inner sep=1pt},
    err/.style={st, draw=red!60!black, thick, fill=red!6},
    fim/.style={st, draw=green!40!black, thick, fill=green!6},
    lbl/.style={font=\scriptsize, text=black!70, inner sep=1pt},
    lble/.style={lbl, text=red!60!black},
    ar/.style={->, draw=black!60},
    are/.style={->, draw=red!60!black}
  ]
  \newcommand{\wf}[3]{{\bfseries #1}\,·\,#2\\[-1pt]{\scriptsize\color{black!65}\ttfamily #3}}
  \def\dx{3.0}\def\dy{1.75}
  \node[st]  (F0) at (0*\dx, 0)      {\wf{0}{preparo}{make pacote}};
  \node[st]  (F1) at (1*\dx, 0)      {\wf{1}{início}{C:\textbackslash exp6}};
  \node[st]  (F2) at (2*\dx, 0)      {\wf{2}{projeto}{abrir exp6.qpf}};
  \node[st]  (F3) at (3*\dx, 0)      {\wf{3}{circuito}{rtl\textbackslash}};
  \node[st]  (F4) at (4*\dx, 0)      {\wf{4}{testbench}{config · cenário}};
  \node[st]  (F5) at (5*\dx, 0)      {\wf{5}{simula}{do simula.do <tb>}};
  \node[err] (F6) at (4*\dx, -\dy)   {\wf{6}{corrige}{ondas · zoom range}};
  \node[st]  (F7) at (5*\dx, -2*\dy) {\wf{7}{todos}{do simula\_todos.do}};
  \node[st]  (F8) at (4*\dx, -2*\dy) {\wf{8}{compila}{Ctrl+L · timing}};
  \node[st]  (F9) at (3*\dx, -2*\dy) {\wf{9}{placa}{Programmer}};
  \node[st]  (FA) at (2*\dx, -2*\dy) {\wf{A}{entrega}{Archive · zip}};
  \node[fim] (FF) at (1*\dx, -2*\dy) {\wf{F}{fim}{pendrive}};

  \draw[ar] (F0) -- (F1);
  \draw[ar] (F1) -- (F2);
  \draw[ar] (F2) -- (F3);
  \draw[ar] (F3) -- (F4);
  \draw[ar] (F4) -- (F5);
  \draw[ar] (F2.north) to[bend left=22] node[lbl, above] {só cenário} (F4.north);
  \draw[ar] (F5) -- (F7);
  \draw[ar] (F7) -- (F8);
  \draw[ar] (F8) -- (F9);
  \draw[ar] (F9) -- (FA);
  \draw[ar] (FA) -- (FF);
  \draw[are] (F5.south west) -- node[lble, pos=0.5, right=3pt] {FALHA, erro} (F6.north east);
  \draw[are] (F7.north west) -- node[lble, pos=0.5, right=3pt] {FALHA} (F6.south east);
  \draw[are] (F8) -- node[lble, right] {erro} (F6);
  \draw[are] (F9.north east) -- node[lble, pos=0.5, left=3pt] {placa $\neq$ simulação} (F6.south west);
  \draw[ar] ([xshift=3mm]F6.north) -- node[lbl, right] {cenário} ([xshift=3mm]F4.south);
  \draw[ar] (F6.north west) -- node[lbl, pos=0.5, left=3pt] {circuito} (F3.south east);
  \draw[ar] (FF.north) -- ++(0,0.6) -| node[lbl, pos=0.25, above] {próximo pedido} (F3.south);
\end{tikzpicture}
\end{center}

\begin{tabularx}{\linewidth}{@{}>{\raggedright\arraybackslash}p{0.13\linewidth} Y >{\raggedright\arraybackslash}p{0.22\linewidth}@{}}
\toprule
Estado & O que fazer & Vai para \\
\midrule
0 preparo & no notebook: \s{make pacote} $\to$ \s{entrega/exp6\_pasta\_t1ba5.zip} no pendrive (projeto + testbenches + cola). Reserva: \s{make entrega} (\s{.qar} e zip do ModelSim) & 1 \\
1 início & extrair numa pasta sem espaço nem acento, ex.: \s{C:\textbackslash exp6}; Quartus: Help $\to$ About mostra a versão & 2 \\
2 projeto & Quartus: File $\to$ Open Project $\to$ \s{quartus\textbackslash de0\_cv\textbackslash exp6.qpf}. Se não abrir (versão mais antiga): Project $\to$ Restore Archived Project com o \s{.qar}, ou New Project Wizard (Cyclone V 5CEBA4F23C7, top \s{sonar2}) + arquivos de \s{rtl} + Assignments $\to$ Import Assignments (\s{exp6.qsf}) & 3 se muda o circuito; 4 se é só cenário \\
3 circuito & editar \s{rtl\textbackslash} nos pontos \s{EXTENSAO} (editor do Quartus ou Notepad++); arquivo novo: Project $\to$ Add/Remove Files in Project & 4 \\
4 testbench & circuito mudou: \s{sonar2\_tb\_config.vh}; cenário: copiar \s{tb\textbackslash cenario\_m\_tb.v} e trocar o nome do módulo & 5 \\
5 simula & ModelSim: File $\to$ Change Directory $\to$ \s{C:\textbackslash exp6}; no Transcript: \s{do simula.do cenario\_x\_tb} & 7 se ``todos os casos OK''; 6 se FALHA ou erro \\
6 corrige & ondas já organizadas pelo \s{wave.do}; colar o \s{wave zoom range} do relatório; editar e repetir \s{do simula.do cenario\_x\_tb} & 3 se o erro é do circuito; 4 se é do cenário \\
7 todos & \s{do simula\_todos.do} (resumo no fim do Transcript) & 8 se 0 com falha; 6 se não \\
8 compila & Quartus: Processing $\to$ Start Compilation (Ctrl+L); Compilation Report $\to$ Timing Analyzer: setup e hold sem slack negativo & 9 se OK; 6 se erro \\
9 placa & Tools $\to$ Programmer $\to$ Hardware Setup: USB-Blaster; Mode JTAG; \s{output\_files\textbackslash exp6.sof}; Start & A se funciona; 6 se a placa difere da simulação \\
A entrega & Project $\to$ Archive Project (\s{exp6\_t1ba5.qar}); zip com \s{rtl}, \s{tb}, \s{*.vh} e \s{*.do}; testar: extrair numa pasta nova e \s{do simula\_todos.do} & F se OK; 6 se não \\
F fim & pasta de volta no pendrive (o notebook não vê o que mudou aqui) & 3 no próximo pedido \\
\bottomrule
\end{tabularx}

\clearpage
% ==========================================================================
\section{Testbench: como modificar (prova)}

\begin{tabularx}{\linewidth}{@{}>{\raggedright\arraybackslash}p{0.21\linewidth} Y >{\raggedright\arraybackslash}p{0.20\linewidth}@{}}
\toprule
Arquivo & O que tem & Mexer quando \\
\midrule
\s{tb/cenario\_x\_tb.v} & o roteiro: casos, comandos, ciclos e conferências, em sequência & sempre (cenário novo) \\
\s{sonar2\_tb\_config.vh} & tudo o que depende do circuito: DUT (portas e parâmetros), tempos da simulação, códigos dos modos, pulsos \s{recebeu\_*}, mensagem esperada, regra do \s{ANDA}, onde o modo aparece & o circuito mudou (procurar \s{EXTENSAO}) \\
\s{sonar2\_tb\_base.vh} & clock, serial (lê e escreve), sensor, conferências, log e relatório da troca de modo & nunca \\
\s{wave.do} & ondas de todos os cenários; nomes de estados e modos lidos da UC; pulsos \s{s\_recebeu\_*} entram sozinhos & sinal novo nas ondas (fim do arquivo) \\
\bottomrule
\end{tabularx}

Estado ou modo novo na UC não muda nada no testbench nem nas ondas: os nomes saem
das linhas \s{parameter nome = 4'hX;} (estado) e \s{parameter NOME = 2'dN;} (modo)
da UC. No cenário, um estado pelo nome: \s{espera\_estado(dut.U1\_UC.muda\_posicao)}.

\begin{minipage}[t]{0.45\linewidth}
\vspace{0pt}
\textbf{Esqueleto} (\s{cp tb/cenario\_m\_tb.v tb/cenario\_x\_tb.v} e trocar o nome do módulo)
\begin{footnotesize}
\begin{verbatim}
`timescale 1ns/1ns
module cenario_x_tb;
`include "sonar2_tb_base.vh"
  initial begin
    inicia;                  // reset, liga
    novo_caso(1, "localizacao");
    ciclo(0, 15, ANDA);
    novo_caso(2, "'a' na espera");
    envia_comando("a", ATENCAO);
    ciclo(1, 25, PARADO);
    finaliza;                // relatorio
  end
endmodule
\end{verbatim}
\end{footnotesize}
{\raggedright
\s{make check TB=cenario\_x\_tb} (Icarus, 10--40\,s)\\
\s{make sim TB=cenario\_x\_tb} (Questa, ondas)\\
\s{make regressao} (todos; acha o cenário novo sozinho)\par}
\end{minipage}\hfill
\begin{minipage}[t]{0.53\linewidth}
\vspace{0pt}
\textbf{Regra:} cada linha começa quando a anterior termina, e as tasks
\s{espera\_*} param até o evento. Para um comando chegar num momento do ciclo,
escreva a linha dele logo depois da task que espera aquele momento.

\begin{tabularx}{\linewidth}{@{}>{\raggedright\arraybackslash}p{0.25\linewidth} Y@{}}
\toprule
Momento & Linhas do cenário \\
\midrule
na espera & \s{envia\_comando("a", ATENCAO);} \s{ciclo(p, cm, PARADO);} \\
na medida & \s{espera\_trigger;} \s{envia\_comando(...);} \s{ciclo(...);} \\
na transmissão & \s{responde\_echo(cm);} \s{espera\_transmissao;} \s{envia\_comando(...);} \s{fim\_ciclo(p, cm, ...);} \\
num estado & \s{responde\_echo(cm);} \s{espera\_estado(dut.U1\_UC.x);} \s{envia\_comando(...);} \s{fim\_ciclo(...);} \\
com o echo em 1 & \s{fork responde\_echo(cm); begin wait (echo\_in); \#(...); envia\_comando(...); end join} (caso 12 dos exemplos) \\
\bottomrule
\end{tabularx}

Antes da bifurcação (espera, medida, transmissão) o comando já vale no próprio
ciclo. Vários comandos seguidos podem passar da espera de 200\,µs: o \s{ciclo}
responde ao trigger que já veio. Padrões prontos: \s{tb/cenario\_exemplos\_tb.v}.
\end{minipage}

\begin{tabularx}{\linewidth}{@{}>{\raggedright\arraybackslash}p{0.40\linewidth} Y@{}}
\toprule
Task & O que faz \\
\midrule
\s{inicia} · \s{finaliza} & reset e liga · desliga, relatório das trocas de modo e resultado \\
\s{novo\_caso(n, "texto")} & número e texto do caso no log e nas ondas (\s{caso}, \s{caso\_texto}) \\
\s{ciclo(p, cm, depois)} & trigger, echo de \s{cm}, \s{fim\_posicao}; confere \s{aaa,ddd\#} e o servo em \s{depois} \\
\s{depois} & \s{ANDA} (regra da config), \s{PARADO} ou a posição exata (0--7) \\
\s{responde\_echo(cm)} + \s{fim\_ciclo(p, cm, depois)} & o mesmo ciclo em duas metades, para pôr algo no meio \\
\s{medida(p, cm)} & echo e mensagem conferida, sem esperar o fim do ciclo (várias medidas por posição) \\
\s{ciclo\_sem\_echo(p, depois)} & sensor não responde: \s{aaa,000\#} e duração $\geq$ \s{T\_TIMEOUT} \\
\s{envia\_comando(c, modo)} & quadro 7E1 (87\,µs) e confere o modo na UC, no LEDR9 e no HEX4 da pág.\ 011 \\
\s{envia\_comando\_paridade\_errada(c, modo)} & quadro com a paridade invertida; confere que o modo não mudou \\
\s{envia\_char(c, PARIDADE\_OK)} & só o quadro, sem conferir \\
\s{espera\_trigger} · \s{espera\_transmissao} · \s{espera\_mensagem} · \s{espera\_fim\_ciclo} · \s{espera\_estado(e)} & param até o evento da medida atual; voltam na hora se ele já aconteceu \\
\s{liga} · \s{desliga} · \s{aplica\_reset} & \s{ligar} = 1 · \s{ligar} = 0 e espera 100\,µs · pulso de 2\,µs no reset \\
\s{confere\_mensagem(p, cm)} · \s{confere\_posicao(p)} · \s{confere\_modo(m)} · \s{confere\_estado(e)} & conferências avulsas \\
\s{confere\_pagina(sel, 24'hXXXXXX)} & os 6 dígitos de uma página; depois volta à página anterior \\
\s{confere\_sinal("nome", sinal, valor)} & qualquer sinal: LED, GPIO novo, interno (\s{dut.s\_...}) \\
\s{confere\_tempo("nome", t, min, max)} & intervalo em ns; \s{t\_trigger}, \s{t\_trigger\_ant}, \s{t\_fim\_ciclo}, \s{T\_ESPERA}, \s{T\_TIMEOUT}, \s{T\_QUADRO} \\
\s{rastro\_uc = 1;} & imprime cada transição da UC, com o nome do estado \\
\s{relata\_janela("nome", t0, \$time)} · \s{marca("texto")} & relatório de qualquer janela (quadros, \s{recebeu\_*}, modo, estados da UC) com \s{wave zoom range} · instante no log \\
\bottomrule
\end{tabularx}

\noindent\begin{minipage}{\linewidth}
\textbf{O circuito mudou: o que mexer no testbench}\\[2pt]
\begin{tabularx}{\linewidth}{@{}>{\raggedright\arraybackslash}p{0.27\linewidth} Y@{}}
\toprule
Mudança no circuito & No testbench \\
\midrule
Estado novo na UC & nada (nome lido da UC; ondas também) \\
Modo novo (código 2 ou 3) & config: \s{localparam DUPLO = 2;} \\
Comando novo (\s{s\_recebeu\_x}) & config: \s{always @(posedge dut.s\_recebeu\_x) registra\_comando("recebeu\_x");} Ondas: nada \\
Caractere de comando trocado & cenário: \s{envia\_comando("N", ...)} \\
Parâmetro novo no \s{sonar2} & config: \s{localparam X\_SIM} reduzido (2\,s $\to$ 200\,µs: $\div 10^4$) e \s{.X(X\_SIM)} no DUT \\
Porta nova (LED, GPIO) & config: \s{wire x\_out;} e \s{.x(x\_out)} no DUT; cenário: \s{confere\_sinal("x", x\_out, 1);} \\
Movimento diferente & config: \s{posicao\_seguinte}; ou a posição exata no \s{ciclo(p, cm, 6)} \\
Várias medidas por posição & cenário: \s{medida(p, cm1); medida(p, cm2); espera\_fim\_ciclo; confere\_posicao(p+1);} \\
Intervalo diferente & cenário: \s{t0 = t\_fim\_ciclo; espera\_trigger; confere\_tempo("intervalo", t\_trigger - t0, min, max);} \\
Formato da mensagem & config: \s{TAM\_MSG} e \s{mensagem\_esperada} \\
Página nova nos displays & cenário: \s{confere\_pagina(3'b100, 24'h...);} \\
Modo mostrado em outro lugar & config, seção 7 (LED e dígito do display) \\
\bottomrule
\end{tabularx}
\end{minipage}

\begin{itemize}[before=\raggedright]
  \item \textbf{Item ``localize a troca de modo'':} o \s{finaliza} imprime, para
        cada troca, \s{rx\_pronto} (dado e paridade), \s{recebeu\_*}, a troca de
        \s{Matual} e as transições da UC do fim do último ciclo no modo antigo até o
        1º \s{dispara\_medida} no novo, mais \s{wave zoom range \ldots} para colar
        no Questa. O \s{wave.do} abre com o grupo Troca\_de\_modo no topo.
  \item \textbf{Icarus x Questa:} no \s{make check}, \s{gerador\_pulso\_tb} e
        \s{interface\_hcsr04\_tb} falham por um artefato do Icarus: o
        \s{gerador\_pulso} faz \s{pulso = 0} e depois \s{1} no mesmo
        \s{always @(*)}, e o trigger ganha bordas de largura zero a cada clock. No
        Questa (\s{make simb}, \s{make regressao}) os dois passam. Nos cenários,
        esperar o trigger com \s{espera\_trigger}, nunca com \s{@(posedge trigger\_out)}.
  \item \textbf{Cenário com erro de sintaxe} quebra o \s{make check},
        \s{make simb} e \s{make sim} de \emph{todos} (eles compilam \s{tb/}
        inteiro): corrigir ou tirar de \s{tb/}. O \s{make regressao} e o
        \s{do simula.do <tb>} isolam o arquivo e rodam os outros.
  \item \textbf{Tempos na simulação:} espera 200\,µs, timeout 20\,ms, quadro
        87\,µs, echo 400\,µs depois do trigger, \s{cm} $\geq$ 1. O \s{sonar2\_tb}
        leva cerca de 1\,min no Icarus.
  \item \textbf{Entrega:} \s{make entrega-teste} monta o zip do ModelSim, extrai
        numa pasta temporária e roda todos os testbenches, como quem corrige.
\end{itemize}

% ==========================================================================
\section{Guia rápido de modificações (circuito)}

Regra: o FD só detecta e executa; a UC decide.
\s{grep -n EXTENSAO rtl/*.v *.vh wave.do} lista todos os pontos de mudança. Depois de
mudar: \s{make check TB=...}, \s{make regressao}, \s{make}, \s{make timing}; arquivo
novo em \s{rtl/}: \s{make sync}.

\begin{tabularx}{\linewidth}{@{}>{\raggedright\arraybackslash}p{0.22\linewidth} >{\raggedright\arraybackslash}p{0.12\linewidth} Y@{}}
\toprule
Pedido do desafio & Onde & Como \\
\midrule
Trocar os caracteres dos comandos & \s{sonar2.v} & \s{CHAR\_ATENCAO = "N"}, \s{CHAR\_VOLTAR = "A"}; no cenário, \s{envia\_comando("N", ...)} \\
Comando novo que muda o modo (ex.: \s{d}) & topo, FD, UC, tb & topo: parâmetro \s{CHAR\_DUPLO}, fio \s{s\_recebeu\_duplo}, ligar na UC e no FD (parâmetro e porta). FD: parâmetro, porta e um comparador igual aos outros. UC: porta, \s{parameter DUPLO = 2'd2;}, transições no \s{case (Matual)}, ramo na bifurcação. tb: \s{localparam DUPLO = 2;} e o \s{always} do pulso \\
Estado novo & \s{sonar2\_uc.v} & código livre B--E; transição no \s{case}; saída no bloco de Moore; estado final de um modo entra no \s{fim\_posicao}. tb e ondas: nada \\
Várias medidas por posição & UC & flag na UC que marca a 1ª medida; na bifurcação: 1ª medida $\to$ estado novo $\to$ \s{espera}; 2ª $\to$ \s{muda\_posicao} (flag volta a 0) \\
Intervalo diferente num modo (ex.: 1\,s) & FD + UC & cópia do U8 com parâmetro novo (o \s{\$clog2} acerta os bits) e \s{conta\_*} próprio vindo da UC; estado de espera só do modo \\
Varredura vai-e-volta & só FD & U4 módulo 14; \s{s\_posicao = (s\_passo < 8) ? s\_passo[2:0] : 14 - s\_passo}; tb: posição exata no \s{ciclo} \\
Ação única (ex.: \s{z} leva o servo a 000) & FD + UC & comparador \s{recebeu\_z}; FF de pedido na UC; na bifurcação, estado que ativa \s{zera\_posicao} (também em \s{preparacao}) \\
Sinal novo nos displays & \s{sonar2.v} & \s{s\_pagina4..7}; sinal interno do FD vira saída \s{db\_...} \\
LED ou GPIO novo & \s{sonar2.v} + \s{.qsf} & porta no topo, pino no Pin Planner (\s{make gui}); tb: porta no DUT da config \\
Entrada assíncrona nova & \s{sonar2.v} & uma instância do \s{sincronizador} antes de ligar na UC ou no FD \\
Tratar ``sem eco'' (807\,cm) & FD & comparador na distância ($>$ 400\,cm) que substitui o valor enviado \\
\bottomrule
\end{tabularx}

O modo já tem 2 bits em todo o caminho (UC $\to$ topo $\to$ HEX4 da pág.\ 011;
LEDR9 = modo $\neq$ 0): um modo novo não muda a largura de nada. Os contadores de
tempo usam \s{N = \$clog2(M)}: qualquer tempo cabe (com \s{N} fixo, um \s{M} maior
que $2^N$ travava a placa e passava na simulação).

\textbf{Ensaio (08/10):} o modo duplo (\s{d}, 2 medidas por posição) feito só nos
pontos \s{EXTENSAO}: 15 edições no RTL, 2 linhas na config e um cenário de
20 linhas; passou de primeira, com os nomes do estado e do modo novos no log e
nas ondas.

% ==========================================================================
\section{Decisões de projeto (perguntas prováveis)}

\begin{tabularx}{\linewidth}{@{}>{\raggedright\arraybackslash}p{0.27\linewidth} Y@{}}
\toprule
Pergunta & Resposta curta \\
\midrule
Por que a troca de modo não perde comandos? & A FSM de modo (\s{Matual}) é atualizada pelo pulso \s{recebeu\_*} em qualquer estado, até no meio da medida ou da transmissão. \\
Por que a bifurcação fica no fim do ciclo? & A única diferença entre os modos é o reposicionamento, que já era o último passo. O ciclo em andamento termina de medir e transmitir. \\
Onde o servo para no modo atenção? & Na UC: o modo 1 pula \s{muda\_posicao}. Não há porta lógica sobre o \s{pwm} nem sobre sinais de controle. \\
Por que conferir a paridade? & Um quadro corrompido poderia virar \s{a} ou \s{v}; quadros com paridade errada são descartados. \\
Por que \s{fim\_posicao} pulsa também em atenção? & Marca o fim de cada ciclo; o tb e o Analog Discovery usam o pulso para sincronizar. \\
O que acontece ao religar? & \s{inicial} força o modo 0 (como o reset) e \s{preparacao} zera a posição. \\
Por que as entradas assíncronas passam por 2 flip-flops? & O Quartus codifica as FSMs em one-hot sem safe state; uma entrada que muda perto da borda do clock pode ser vista diferente por FFs diferentes e deixar a FSM num estado ilegal (HEX5 parado em 3 na bancada, 06/10). O \s{sincronizador} entrega um valor único por clock e dá tempo para a metaestabilidade se resolver; vale para \s{ligar}, \s{echo} e \s{entrada\_serial}. \\
O atraso de 40\,ns atrapalha? & Não: o echo atrasa nas duas bordas (a largura não muda) e 40\,ns são 0,5\,\% de um bit da serial. \\
Por que o modo tem 2 bits se há 2 modos? & Reserva para a prova: um modo novo só muda a UC. O LEDR9 acende fora da localização e o HEX4 da pág.\ 011 mostra o número. \\
Por que 807\,cm? & É o pulso de cerca de 47\,ms que o sensor gera sem eco; o circuito mede corretamente. \\
Recursos do projeto & 259 ALMs (1\%), 331 registradores, 64 pinos; pior setup 15,14\,ns, pior hold 0,149\,ns; Fmax 214\,MHz. \\
\bottomrule
\end{tabularx}

\end{document}
